% ============================================================================
% Hindi UI and structural strings for the One Physics Book series.
% Loaded by styles/onephysics.sty when \booklang is hi.
% UTF-8 Devanagari throughout. Build with XeLaTeX (see *_hi.tex entries).
% ============================================================================

% Theorem / statement environment titles
\newcommand{\omnameDefinition}{परिभाषा}
\newcommand{\omnameTheorem}{प्रमेय}
\newcommand{\omnameProposition}{प्रस्ताव}
\newcommand{\omnameLemma}{उपप्रमेय}
\newcommand{\omnameCorollary}{उपप्रमेयिका}
\newcommand{\omnameMethod}{विधि}
\newcommand{\omnameExample}{उदाहरण}
\newcommand{\omnameNotation}{संकेतन}
\newcommand{\omnameRemark}{टिप्पणी}
\newcommand{\omnameExercise}{अभ्यास}
\newcommand{\omnameExercises}{अभ्यास}
\newcommand{\omnameProblem}{समस्या}
\newcommand{\omnameProblems}{समस्याएँ}
\newcommand{\omnameProof}{उपपत्ति}

% Admitted results and solutions appendix headers
\newcommand{\omadmittedtext}{इस स्तर पर स्वीकृत।}
\newcommand{\omsolutionof}{हल ---}
\newcommand{\ompage}{पृष्ठ}
\newcommand{\omnameGotoSolution}{हल पृ.}

% Misc text macros
\newcommand{\st}{\text{ ऐसा कि }}

% Cover / title page
\newcommand{\bookmaintitle}{One Physics Book}
\newcommand{\booktagline}{सब पर राज करने वाली एक भौतिकी की पुस्तक}
\newcommand{\bookdescription}{किंडरगार्टन से स्नातक की डिग्री के अंत तक भौतिकी}
\newcommand{\bookcontributors}{One Physics Book के योगदानकर्ता}
\newcommand{\bookversionlabel}{संस्करण}

% Colophon
\newcommand{\omcolophonsource}{स्रोत कोड, समस्याएँ और योगदान:}

% Solutions appendix part title
\newcommand{\omsolutionspart}{अभ्यासों के हल}

% Part titles (Grades 1–12; parallel to English Grade N)
\@namedef{omstr@part.grade1}{कक्षा 1 (आयु 6--7)}
\@namedef{omstr@part.grade2}{कक्षा 2 (आयु 7--8)}
\@namedef{omstr@part.grade3}{कक्षा 3 (आयु 8--9)}
\@namedef{omstr@part.grade4}{कक्षा 4 (आयु 9--10)}
\@namedef{omstr@part.grade5}{कक्षा 5 (आयु 10--11)}
\@namedef{omstr@part.grade6}{कक्षा 6 (आयु 11--12)}
\@namedef{omstr@part.grade7}{कक्षा 7 (आयु 12--13)}
\@namedef{omstr@part.grade8}{कक्षा 8 (आयु 13--14)}
\@namedef{omstr@part.grade9}{कक्षा 9 (आयु 14--15)}
\@namedef{omstr@part.grade10}{कक्षा 10 (आयु 15--16)}
\@namedef{omstr@part.grade11}{कक्षा 11 (आयु 16--17)}
\@namedef{omstr@part.grade12}{कक्षा 12 (आयु 17--18)}

% Part titles (Bachelor Years 1–3)
\@namedef{omstr@part.bachelor1}{स्नातक वर्ष 1 (आयु 18--19)}
\@namedef{omstr@part.bachelor2}{स्नातक वर्ष 2 (आयु 19--20)}
\@namedef{omstr@part.bachelor3}{स्नातक वर्ष 3 (आयु 20--21)}

% Plural names + \omnameChapter, used by the \crefname declarations in
% onephysics.sty.
\newcommand{\omnameDefinitions}{परिभाषाएँ}
\newcommand{\omnameTheorems}{प्रमेय}
\newcommand{\omnamePropositions}{प्रस्ताव}
\newcommand{\omnameLemmas}{उपप्रमेय}
\newcommand{\omnameCorollarys}{उपप्रमेयिकाएँ}
\newcommand{\omnameMethods}{विधियाँ}
\newcommand{\omnameExamples}{उदाहरण}
\newcommand{\omnameNotations}{संकेतन}
\newcommand{\omnameRemarks}{टिप्पणियाँ}
\newcommand{\omnameChapter}{अध्याय}
\newcommand{\omnameChapters}{अध्याय}

% LaTeX's own structural names (babel is not used for Hindi).
\renewcommand{\chaptername}{अध्याय}
\renewcommand{\partname}{भाग}
\renewcommand{\contentsname}{विषय-सूची}
\renewcommand{\appendixname}{परिशिष्ट}
\renewcommand{\indexname}{अनुक्रमणिका}
